\section{\sc Internships}
% Location and dates sit on the heading lines, so the bullets use the full width
\textbf{Tata Research Design and Development Center} \hfill Pune, India\\
{\parfillskip=0pt \emph{Research Intern in Machine Learning, \href{https://m-raj.github.io/pages/projects/particle_classifier_optimisation.html}{Project Website}}\hfill Winter, 2020\par}
\begin{list1}
    \item Developed a \textbf{machine-learning-driven framework for multi-objective optimization} of product designs, integrating CFD simulations with surrogate modeling.

\item Built a \textbf{machine learning surrogate} model to rapidly predict separation factor and pressure drop, enabling efficient exploration of the particle-classifier design space.
\item Performed \textbf{multi-objective optimization} to identify Pareto-optimal designs balancing separation performance against pressure drop.
\end{list1}

\textbf{RWTH Aachen University} \hfill Aachen, Germany\\
{\parfillskip=0pt \emph{Research Intern in Computational Mechanics, \href{https://m-raj.github.io/projects/nacre-composites/}{Project Website}}\hfill Summer, 2019\par}
\begin{list1}
    % \item Modeled nacre microstructure using and conducted finite-element-method (FEM) simulations to investigate the effect of material composition.
    \item Modeled the 3D brick-and-mortar microstructure of nacre using HyperMesh, representing the composite architecture of brittle calcium carbonate and ductile organic matrix.

\item Conducted LS-DYNA simulations to characterize the mechanical properties and investigate the effect of material composition on the toughness of nacre-based composites.

\item Demonstrated that increasing the organic-matrix fraction enhances composite toughness but results in a corresponding reduction in overall strength.
\end{list1}
\endinput
